\section{Zadanie 1}
% opis 
Porównujemy różne algorytmy mnożące macierze 
\textbf{Działanie mm}:
\begin{enumerate}
    \item pętla po indeksach wierszy macierzy A
    \item pętla po indeksach kulmn macierzy B
    \item pętla po indekasch kolumn macierzy A 
    \item zliczanie każdego elementu po kolei (mnożąc odpowiednie miejsca w wierszach i kolumnach
\end{enumerate}
\textbf{Działanie mm1}:
\begin{enumerate}
    \item pętla po indeksach wierszy macierzy A
    \item pętla po indeksach kulmn macierzy B
    \item mnożenie macierzowe odpowiednich wieszych A, kolumn B
\end{enumerate}
\textbf{Domyślne mnożenie macierzy}:\\
Nie mogłem znaleźć jak to jest skonstruowane, najprawdopodobniej pętle pętle sa inaczej skonstruowane.\\

\begin{center}
\resizebox{0.6 \textwidth}{!}{%
\begin{tabular}{|l|l|l|l|}
\hline
$n$   & mm         & mm1         & default         \\ \hline
100   & $5.28124$s & $0.052098$s & $0.000175953$   \\ \hline
500   & $671.192$s & $1.42004$s  & $0.00494003$   \\ \hline
\end{tabular}%
}
\end{center}

% wnioski 

\noindent 
\textbf{Wnioski}:\\ widzimy że dodawanie kolejnej pętli znacznie zwiększa złożoność obliczeniową. W pierwszym przypadku mamy złożoność $n^3$, w drugim $n^2$. Domyślny algorytm musi korzystać z innych sposób, widzimy że roobi to znacznie szyciej niż nasza implementacja. 